El hardware utilizado para realizar los experimentos corresponde a un procesador Intel Core i5-12450H, 16 GB de memoria RAM DDR4 a 3200 MT/s y sistema operativo Windows 11. Los algoritmos fueron ejecutados en un entorno Ubuntu mediante WSL2, el cual tenia asignado aproximadamente 7,6 GiB de memoria RAM. Para la compilación de los programas en C++ se utilizó g++ (GCC) versión 13.3.0, utilizando el estándar C++17.

Los datasets que se usaron corresponden a los archivos generados mediante array\_generator.py y matrix\_generator.py, siguiendo las instrucciones y configuraciones definidas para la tarea. Para los algoritmos de ordenamiento se utilizaron arreglos con distribuciones de tipo ascendente, descendente y aleatoria, considerando los dominios D1 y D7 y tres muestras por configuración. Además de los tamaños definidos para las pruebas, se incorporaron tamaños intermedios y se quitaron otros, esto ya que con algunos de los tamaños entregados por la tarea, los tiempos de ejecucion eran superiores a 10 minutos por arreglo, por lo que el tamaño máximo de un arreglo para los algoritmos de ordenamiento es de $10^5$, descartando los arreglos de tamaño de $10^6$ y $10^7$ entregados en la tarea inicialmente. Luego para la multiplicación de matrices se utilizaron matrices cuadradas de dimensiones \(2^4\), \(2^6\), \(2^8\) y \(2^{10}\), de tipo dispersa, diagonal y densa, con dominios D0 y D10 y tres muestras por configuración.